% ЗГЕНЕРОВАНО scripts/build_format_v2_chapter.py — не редагувати вручну.
% Джерело: editorial/cards/5.*.json, format-v2/rozdil-5.json, paraleli.json.

\naukchapter{5}{4}{Глава 5. Розчинення}
\naukintro{Пʼята глава — чотири вірші й один рефрен. Рефрен тут не прикраса: він стоїть у кінці кожної строфи тими самими словами, і саме повтор і є думкою. Глава не описує розчинення, а показує, що робити з ним, коли відкидати вже нічого: у першому вірші це сказано прямо — у чистого немає звʼязку ні з чим, тож і відкидати йому нічого.}
\vsec{5.1}
\vdeva{न ते सङ्गो ऽस्ति केनापि किं शुद्धस् त्यक्तुम् इच्छसि । सङ्घातविलयं कुर्वन्न् एवम् एव लयं व्रज ॥}
\vtranslit{на те санґо ’сті кенапі кім шуддгас тяктум іччасі | санґгатавілаям курванн евам ева лаям враджа ||}
\vperek{У тебе немає зв'язку ні з чим — то чого ж ти, чистий, хочеш відкинути? Здійснюючи розчинення скупчення, так само йди до розчинення.}
\vcom{Построкове зіставлення: редакторська чернетка, не рецензовано.}
\vcom{Рядок а}
\vdeva{न ते सङ्गो ऽस्ति केनापि किं शुद्धस् त्यक्तुम् इच्छसि}
\vtranslit{на те санґо ’сті кенапі кім шуддгас тяктум іччасі}
\vperek{У тебе немає зв'язку ні з чим — то чого ж ти, чистий, хочеш відкинути?}
\vcom{na ... saṅgo 'sti заперечує прив'язаність; kenāpi вказує на відсутність зв'язку з будь-чим, kiṃ ... icchasi формує питання, tyaktum — інфінітив «відкинути». Тлумачення стосується внутрішньої свободи, а не наказу полишити родину чи працю.}
\vcom{Рядок б}
\vdeva{सङ्घातविलयं कुर्वन्न् एवम् एव लयं व्रज}
\vtranslit{санґгатавілаям курванн евам ева лаям враджа}
\vperek{Здійснюючи розчинення скупчення, так само йди до розчинення.}
\vcom{saṅghātavilayaṃ — розчинення сукупності, kurvann — дієприкметник «здійснюючи», vraja — наказова форма «йди». evam eva layaṃ vraja — рефрен усіх чотирьох віршів глави 5; розрізнення vilaya / laya як філософських термінів потребує незалежної фахової звірки.}
\vcom{Питання kim ... icchasi («чого хочеш») викриває безглуздість самого наміру відректися: відкидають те, до чого прив'язані, а зв'язку немає (na te saṅgaḥ asti kenāpi). Слово saṅghāta («скупчення») — технічне: тіло описане не як річ, а як агрегат, суму. Тому vilaya («розчинення») і laya («розчинення») стоять в одному вірші двічі: перше — те, що відбувається з тілом само собою, друге — те, до чого слід іти. Різниця в тому, що перше трапляється, а друге обирають. РЕФРЕН ГЛАВИ. Усі чотири вірші глави 5 закінчуються однаково: evam eva layaṃ vraja («так само йди до розчинення»). Це не чотири окремі повчання, а одна літанія з незмінним приспівом: змінюється тільки перша половина — обґрунтування, а висновок лишається той самий. Тому главу слід читати як одне ціле.}
\vpar{Українська паралель (наша інтерпретація)}{Григорій Сковорода, «Асхань. Симфонія, наречена книга Асхань, про пізнання самого себе».}{Наша паралель: у Сковороди пізнання себе теж нічого не додає й нічого не відкидає — воно знімає перешкоду, і тільки. Наша інтерпретація, не твердження про вплив.}
\vfig{ilustratsii/5-1-neti.pdf}{Відкидати нічого: у чистого немає звʼязку}{Схема наша, не історична ілюстрація.}

\vsec{5.2}
\vdeva{उदेति भवतो विश्वं वारिधेर् इव बुद्बुदः । इति ज्ञात्वा एकम् आत्मानम् एवम् एव लयं व्रज ॥}
\vtranslit{удеті бгавато вішвам варідгер іва будбудах | іті джнятва екам атманам евам ева лаям враджа ||}
\vperek{Світ постає з тебе, як бульбашка з океану. Знаючи, що Атман єдиний, — так само йди до розчинення.}
\vcom{Построкове зіставлення: редакторська чернетка, не рецензовано.}
\vcom{Рядок а}
\vdeva{उदेति भवतो विश्वं वारिधेर् इव बुद्बुदः}
\vtranslit{удеті бгавато вішвам варідгер іва будбудах}
\vperek{Світ постає з тебе, як бульбашка з океану.}
\vcom{udeti означає «постає»; bhavataḥ допускає значення «з тебе» або «твого», vāridher iva budbudaḥ порівнює світ із бульбашкою в океані. Це образ залежного виникнення, а не самостійний доказ неіснування світу.}
\vcom{Рядок б}
\vdeva{इति ज्ञात्वा एकम् आत्मानम् एवम् एव लयं व्रज}
\vtranslit{іті джнятва екам атманам евам ева лаям враджа}
\vperek{Знаючи, що Атман єдиний, — так само йди до розчинення.}
\vcom{iti jñātvā — «так знаючи», ekam ātmānam стосується єдності Атмана. evam eva layaṃ vraja — наказовий рефрен усіх чотирьох віршів цієї глави; український відповідник збережено без змін.}
\vcom{Порівняння з бульбашкою точніше за звичне «хвиля в морі» (2.10, 2.23–2.25): хвиля є форма тієї ж води, а бульбашка — порожнеча, обтягнута плівкою з тієї ж води. Вона не має власного вмісту. bhavataḥ стоїть в ablativus/genitivus, і обидва прочитання працюють: «з тебе» (джерело) і «твого» (належність). Висновок вірша не «світ нереальний», а «Атман єдиний» — і саме з єдності випливає розчинення, а не навпаки. РЕФРЕН ГЛАВИ. Усі чотири вірші глави 5 закінчуються однаково: evam eva layaṃ vraja («так само йди до розчинення»). Це не чотири окремі повчання, а одна літанія з незмінним приспівом: змінюється тільки перша половина — обґрунтування, а висновок лишається той самий. Тому главу слід читати як одне ціле.}
\vpar{Українська паралель (наша інтерпретація)}{Богдан-Ігор Антонич, «Ротації» (1938).}{Наша паралель: в Антонича світ теж постає з того, хто на нього дивиться, як бульбашка з океану. Наше порівняння — за образом.}
\vfig{ilustratsii/5-2-space.pdf}{Бульбашка й океан: світ постає з тебе}{Схема наша, не історична ілюстрація.}
\vfig{ilustratsii-naratyvni/5-2-bulbashka-z-okeanu.pdf}{Людина сидить на березі, а з океану одна за одною піднімаються бульбашки}{Наша ілюстрація до вірша 5.2 — не історичне зображення.}

\vsec{5.3}
\vdeva{प्रत्यक्षम् अप्य् अवस्तुत्वाद् विश्वं नास्त्य् अमले त्वयि । रज्जुसर्प इव व्यक्तम् एवम् एव लयं व्रज ॥}
\vtranslit{пратякшам апй авастутвад вішвам настй амале тваї | раджджусарпа іва в'яктам евам ева лаям враджа ||}
\vperek{Хоч світ і перед очима, він не існує в тобі, незаплямованому, — бо він не річ, як змій у мотузці, явно. Так само йди до розчинення.}
\vcom{Построкове зіставлення: редакторська чернетка, не рецензовано.}
\vcom{Рядок а}
\vdeva{प्रत्यक्षम् अप्य् अवस्तुत्वाद् विश्वं नास्त्य् अमले त्वयि}
\vtranslit{пратякшам апй авастутвад вішвам настй амале тваї}
\vperek{Хоч світ і перед очима, він не існує в тобі, незаплямованому, — бо він не річ,}
\vcom{pratyakṣam api означає «хоч і видимий перед очима»; avastutvād — «через відсутність самостійної речовості/сутності», viśvaṃ nāsti — «світ не існує»; amale tvayi звертається до незаплямованого адресата. Це відтворення редакторської картки, а не нове критичне читання.}
\vcom{Рядок б}
\vdeva{रज्जुसर्प इव व्यक्तम् एवम् एव लयं व्रज}
\vtranslit{раджджусарпа іва в'яктам евам ева лаям враджа}
\vperek{як змій у мотузці, явно. Так само йди до розчинення.}
\vcom{rajjusarpa iva — порівняння з уявним змієм у мотузці, vyaktam — «виразно, явно». evam eva layaṃ vraja зберігає незмінний рефрен глави; філософське пояснення образу слід відрізняти від лексичного значення санскриту.}
\vcom{Найсильніше тут — avastutvāt («через несуттєвість»). Це не «світ оманливий», а точніше: світ не є vastu, річчю, самостійною сутністю. pratyakṣam api («хоч і перед очима») визнає видимість повністю — вірш не заперечує, що ми бачимо. Заперечується лише те, що видиме має власне буття. Змій у мотузці (той самий приклад, що в 2.9) поданий тут з vyaktam («явно»): помилка не тьмяна, вона гранично чітка — і саме тому переконлива. РЕФРЕН ГЛАВИ. Усі чотири вірші глави 5 закінчуються однаково: evam eva layaṃ vraja («так само йди до розчинення»). Це не чотири окремі повчання, а одна літанія з незмінним приспівом: змінюється тільки перша половина — обґрунтування, а висновок лишається той самий. Тому главу слід читати як одне ціле.}
\vpar{Українська паралель (наша інтерпретація)}{Іван Франко, «Сойчине крило» (1905).}{Наша паралель: у Франка видиме й справжнє теж не збігаються, і саме на цьому розходженні тримається оповідь. Наше порівняння — типологічне, не про джерело.}
\vfig{ilustratsii/5-3-illusion.pdf}{Світ перед очима й змій у мотузці}{Схема наша, не історична ілюстрація.}

\vsec{5.4}
\vdeva{समदुःखसुखः पूर्ण आशानैराश्ययोः समः । समजीवितमृत्युः सन्न् एवम् एव लयं व्रज ॥}
\vtranslit{самадухкхасукхах пурна ашанайрашяйох самах | самаджівітамрітюх санн евам ева лаям враджа ||}
\vperek{Рівний у стражданні й радості, повний, рівний у сподіванні й безнадії, рівний у житті й смерті, — так само йди до розчинення.}
\vcom{Построкове зіставлення: редакторська чернетка, не рецензовано.}
\vcom{Рядок а}
\vdeva{समदुःखसुखः पूर्ण आशानैराश्ययोः समः}
\vtranslit{самадухкхасукхах пурна ашанайрашяйох самах}
\vperek{Рівний у стражданні й радості, повний, рівний у сподіванні й безнадії,}
\vcom{samaduḥkhasukhaḥ — рівність до страждання й радості; pūrṇaḥ — «повний»; āśānairāśyayoḥ samaḥ — рівний у парі сподівання/безнадія. Два протиставлення містяться в одному друкованому піврядку; це не два самостійні переклади.}
\vcom{Рядок б}
\vdeva{समजीवितमृत्युः सन्न् एवम् एव लयं व्रज}
\vtranslit{самаджівітамрітюх санн евам ева лаям враджа}
\vperek{рівний у житті й смерті, — так само йди до розчинення.}
\vcom{samajīvitamṛtyuḥ — рівність у житті й смерті, san — дієприкметник «бувши». evam eva layaṃ vraja — той самий наказовий рефрен, що в 5.1–5.3; український порядок слів залишено без змін.}
\vcom{Вірш складається з трьох пар, і всі три побудовані на sama («рівний»): страждання/радість, сподівання/безнадія, життя/смерть. Пари охоплюють три виміри — відчуття, очікування, буття. Слово pūrṇaḥ («повний») стоїть між першою і другою парою не випадково: рівність не порожнеча, а повнота, якій не бракує нічого з обох боків. Це єдиний вірш глави, де немає ні прикладу, ні заперечення — лише перелік ознак. Тим самим літанія закінчується не аргументом, а описом того, хто йде. РЕФРЕН ГЛАВИ. Усі чотири вірші глави 5 закінчуються однаково: evam eva layaṃ vraja («так само йди до розчинення»). Це не чотири окремі повчання, а одна літанія з незмінним приспівом: змінюється тільки перша половина — обґрунтування, а висновок лишається той самий. Тому главу слід читати як одне ціле.}
\vpar{Українська паралель (наша інтерпретація)}{Леся Українка, «Камінний господар» (1912).}{Наша паралель: у Лесі Українки рівність перед останнім теж не залежить від надії чи її відсутності. Наше порівняння — за будовою думки.}
\vfig{ilustratsii/5-4-balance.pdf}{Дві шальки, що стоять рівно: життя й смерть}{Схема наша, не історична ілюстрація.}

